\section{Related Work}
\label{sec:related_work}

\para{Number representation in language models.}
A growing body of work investigates how language models encode numerical information. \citet{wallace2019nlp} showed that even early embeddings (GloVe, ELMo, BERT) encode numerical magnitude, though models cannot extrapolate beyond their training range. More recently, \citet{levy2024language} demonstrated that LLMs encode numbers using \emph{per-digit circular features in base 10}: circular probes targeting $\cos(2\pi t / 10)$ and $\sin(2\pi t / 10)$ per digit position achieve 91--100\% accuracy. \citet{kadlcik2025pretrained} confirmed and extended this finding, showing that sinusoidal Fourier-basis probes achieve 94--100\% accuracy across model families. \citet{yuchi2026llms} found that linear probes recover log-magnitudes with $\sim$2.3\% median relative error and that comparison classifiers achieve $>$90\% accuracy from hidden states even when verbalized accuracy is only 50--70\%. Our work builds on these findings by studying how such number representations support \emph{composite} operations like interval membership.

\para{Mechanistic interpretability of arithmetic.}
\citet{hanna2023how} identified the exact circuit GPT-2 uses for greater-than comparison: attention heads in layers~5--9 attend to the reference year, and MLPs~8--11 compute the comparison by producing an upper-triangular pattern in logit space. Critically, they found that GPT-2 has no less-than circuit and overgeneralizes greater-than. \citet{stolfo2023mechanistic} used causal mediation analysis to identify three phases of arithmetic processing: early MLPs encode operands, middle-layer attention transfers information, and late MLPs compute results. \citet{quirke2024understanding} showed that small transformers learn cascading carry/borrow algorithms for addition and subtraction. Our work extends these analyses to the composite operation of interval membership, which requires coordinating two boundary checks.

\para{Counting and modular arithmetic.}
\citet{hasani2026mechanistic} showed that counting information accumulates progressively across layers, with count information localized at boundary tokens in middle-to-late layers. Accuracy degrades sharply beyond $\sim$10 items. \citet{chang2024language} demonstrated that transformers cannot generalize counting inductively without positional embeddings, and \citet{zhang2024transformers} noted that transformers are limited to the TC$^0$ complexity class, constraining deep counting. Our alphabetic ring counting task---which requires modular arithmetic over a 26-element group---extends these findings to a domain where digit-level heuristics cannot apply.

\para{Probing methodology.}
Linear probing, introduced by \citet{alain2017understanding}, is the standard method for testing what information is linearly accessible in hidden states. \citet{hewitt2019designing} highlighted the risk that expressive probes can memorize random labels, motivating the use of shuffled-label controls. We adopt both linear probes and shuffled-label baselines in our experiments.

\para{Behavior--representation dissociation.}
\citet{yuchi2026llms} demonstrated that LLMs ``know more about numbers than they can say,'' finding that internal representations encode numerical information far more accurately than the model's verbalized outputs. Our work extends this finding to interval membership and modular arithmetic, documenting the most extreme case of this dissociation yet observed: 0\% behavioral accuracy with 58\% internal decodability for ring counting.
